A normal mixture with the medians and spreads of the two Danish groups of 2013 is a poor guide: it is far less separated than the data, and its power at $n=6{,}800$ is 9 per cent where the data give $p=0.005$. Power was therefore measured directly, by drawing subsamples of 150, 300 and 700 people from the Danish refined transfer-reliant group of 2013, with probability proportional to the household weight, and applying the unweighted test to each (ten subsamples of each size, run within LIS). The test rejected at the five per cent level in two, five and seven of the ten subsamples. At the Swedish sizes, between 150 and 750, a bimodality as marked as the Danish one would thus be missed in roughly a third to four fifths of samples; only in the two larger Swedish waves, 2007 and 2010, would it more often than not have been found. The Swedish results are therefore weak evidence against it at most.